\begin{table}[t]
  \centering
  {\small
  \setlength{\tabcolsep}{5pt}
  \begin{tabular}{lccc}
    \toprule
    Method & MCAR $\downarrow$ & MAR $\downarrow$ & NMAR $\downarrow$ \\
    \midrule
    \textbf{CAIR\ (Ours)} & \textbf{2.66} & \textbf{10.33} & \textbf{23.50} \\
    \midrule
    linear interp & \underline{2.91} & \underline{12.34} & \underline{28.94} \\
    MICE          & 3.26 & 19.83 & 49.35 \\
    missForest    & 3.33 & 16.54 & 40.90 \\
    hot-deck      & 5.04 & 16.06 & 43.44 \\
    $k$NN         & 5.19 & 15.45 & 43.54 \\
    LOCF          & 6.22 & 19.84 & 34.44 \\
    Fourier       & 9.91 & 18.47 & 32.72 \\
    mean          & 28.59 & 28.81 & 54.86 \\
    GP-VAE        & 26.46 & 40.49 & 70.39 \\
    M-RNN         & 44.39 & 44.34 & 64.63 \\
    \bottomrule
  \end{tabular}}
  \caption{\textbf{Realistic missingness imputation on AI-READI CGM} (all $352$ test
  participants; RMSE mg/dL averaged over six missingness rates, lower is
  better). CAIR\ is best under all three mechanisms and its margin grows with
  difficulty. The two generic learned imputers fall below every non-constant
  baseline throughout. \textbf{Bold} = best, \underline{underline} = second best, per column.}
  \label{tab:mech}
\end{table}

\begin{table}[t]
  \centering
  {\small
  \setlength{\tabcolsep}{3.6pt}
  \begin{tabular}{lccccc}
    \toprule
    Method & 15\,min & 30\,min & 45\,min & 60\,min & \textbf{mean} $\downarrow$ \\
    \midrule
    \textbf{CAIR\ (Ours)}   & \underline{2.84} & \underline{5.03} & \textbf{6.66} & \textbf{8.23} & \textbf{5.69} \\
    akima                    & \textbf{2.76} & \textbf{5.01} & \underline{6.87} & 8.61 & \underline{5.81} \\
    PCHIP                    & 2.84 & 5.14 & 6.92 & \underline{8.58} & 5.87 \\
    linear                   & 3.20 & 5.64 & 7.45 & 9.09 & 6.34 \\
    SAITS                    & 4.60 & 7.54 & 9.76 & 11.79 & 8.42 \\
    BRITS                    & 6.25 & 9.42 & 11.67 & 13.33 & 10.17 \\
    \bottomrule
  \end{tabular}}
  \caption{\textbf{Imputation performance vs. gap length} on AI-READI CGM (all $352$ test
  participants; RMSE mg/dL). CAIR\ is best on the mean and at
  $45$-$60$\,min, and second by a small margin at $15$-$30$\,min.
  \textbf{Bold} = best, \underline{underline} = second, per column.}
  \label{tab:shortgap}
\end{table}

\begin{table}[t]
  \centering
  {\small
  \setlength{\tabcolsep}{3.2pt}
  \begin{tabular}{lcccccc}
    \toprule
    Method & meal & sleep & asc & dip & comb & \textbf{AVG} $\downarrow$ \\
    \midrule
    \textbf{CAIR\ (Ours)} & \textbf{15.25} & \textbf{24.94} & 9.01 & 7.49 & \textbf{16.37} & \textbf{14.61} \\
    \midrule
    PCHIP           & \underline{16.56} & \underline{29.26} & \textbf{7.86} & 7.36 & 18.65 & \underline{15.94} \\
    linear          & 16.97 & 29.27 & 8.42 & \underline{7.31} & \underline{18.30} & 16.05 \\
    akima           & 17.82 & 36.73 & \underline{8.02} & \textbf{7.24} & 19.64 & 17.89 \\
    Kalman          & 17.97 & 34.70 & 8.88 & 8.09 & 20.29 & 17.99 \\
    Savitzky-Golay & 17.27 & 29.45 & 13.41 & 11.22 & 19.43 & 18.15 \\
    SAITS           & 21.27 & 32.02 & 11.22 & 10.22 & 19.70 & 18.89 \\
    BRITS           & 32.72 & 43.61 & 24.25 & 20.81 & 31.45 & 30.57 \\
    \bottomrule
  \end{tabular}}
  \caption{\textbf{Imputation of physiologically-significant missingness} on AI-READI (all $352$
  test participants; RMSE mg/dL, lower is better). CAIR\ attains the best
  average over the full baseline suite, and the gain concentrates in the long
  \emph{sleep} blocks. The four weakest classical fills are in the supplement.
  \textbf{Bold} = best, \underline{underline} = second
  best, per column.}
  \label{tab:phys}
\end{table}

\begin{table}[t]
  \centering
  {\small
  \setlength{\tabcolsep}{4pt}
  \begin{tabular}{@{}lccc@{}}
    \toprule
    \multicolumn{4}{@{}l}{\emph{(a) Physiological vs.\ random-gap training masks}}\\
    Training masks & \multicolumn{3}{c@{}}{ABP RMSE $\downarrow$}\\
    & MCAR & MAR & NMAR \\
    \midrule
    linear                  & \underline{4.40} & 9.23 & \underline{11.85} \\
    CAIR, CGM masks        & 5.47 & \underline{8.74} & 12.10 \\
    \textbf{CAIR, random-gap} & \textbf{4.30} & \textbf{8.60} & \textbf{11.42} \\
    \midrule
    \multicolumn{4}{@{}l}{\emph{(b) Clinical-burden recovery}}\\
    Method & \multicolumn{3}{c@{}}{MRR $\uparrow$}\\
    & MCAR & MAR & NMAR \\
    \midrule
    linear                & $-0.31$ & $-0.30$ & $+0.14$ \\
    $k$NN                 & $\mathbf{+0.79}$ & $\mathbf{+0.74}$ & $\mathbf{+0.65}$ \\
    \textbf{CAIR (ours)} & \underline{$+0.71$} & \underline{$+0.64$} & \underline{$+0.64$} \\
    \bottomrule
  \end{tabular}}
  \caption{\textbf{Two analyses behind the MIMIC-III results} (ABP).
  \textbf{(a)} The \emph{identical} architecture, trained on the same ABP data,
  is less accurate than linear under MCAR and NMAR when trained on masks carried
  over from glucose physiology, and more accurate under all three under the
  random-gap curriculum.
  \textbf{(b)} Linear is \emph{worse than mean-fill} ($\mathrm{MRR}<0$) under
  MCAR/MAR; CAIR\ and $k$NN recover most of the burden.
  \textbf{Bold} = best, \underline{underline} = second, per column.
  A third analysis, in which no imputer separates from any other when
  predicting mortality from arterial pressure alone, is reported in
  Sec.~\ref{sec:exp-down}.}
  \label{tab:analysis}
\end{table}

\begin{table}[t]
\centering
  \setlength{\tabcolsep}{1mm}
  {\small
\begin{tabular}{lcccccccc}
\toprule
& \multicolumn{4}{c}{Heart rate (bpm) $\downarrow$} & \multicolumn{4}{c}{Resp.\ (br/min) $\downarrow$} \\
\cmidrule(lr){2-5}\cmidrule(lr){6-9}
Method & 15 & 30 & 45 & 60 & 15 & 30 & 45 & 60 \\
\midrule
Akima              & 4.54 & 5.93 & 6.86 & 7.87 & 2.73 & 3.23 & 3.72 & 4.30 \\
AR (bidir.)        & 4.71 & 5.62 & \underline{6.09} & \underline{6.72} & 2.77 & 3.09 & \underline{3.33} & \underline{3.57} \\
PCHIP              & 4.45 & 5.62 & 6.27 & 7.05 & 2.63 & 3.07 & 3.43 & 3.76 \\
Linear             & \underline{4.40} & \underline{5.52} & 6.12 & 6.86 & \underline{2.61} & \underline{3.00} & 3.34 & 3.65 \\
\midrule
\multicolumn{9}{l}{CAIR\ (ours)} \\
\quad univariate   & 4.96 & 6.49 & 6.74 & 7.51 & 2.90 & 3.92 & 4.50 & 5.21 \\
\quad \textbf{multivariate} & \textbf{3.50} & \textbf{4.19} & \textbf{4.29} & \textbf{4.43} & \textbf{1.88} & \textbf{2.54} & \textbf{2.62} & \textbf{2.72} \\
\bottomrule
\end{tabular}}
  \caption{\textbf{Cross-modal conditioning is what transfers} (AI-READI,
gap-length protocol, RMSE in native units, lower is better). Column headings are
gap lengths in minutes; ``Resp.'' is respiration and ``AR (bidir.)'' is
bidirectional AR. Univariate
CAIR\ is less accurate than linear at every gap length on both signals; adding
the co-recorded context makes it best everywhere. On glucose the same context
adds nothing (Sec.~\ref{sec:exp-other}). \textbf{Bold} = best, \underline{underline} = second, per column.}
  \label{tab:hr-resp}
\end{table}

\begin{table}[t]
  \centering
  {\small
  \setlength{\tabcolsep}{1mm}
  \begin{tabular}{lcccccc}
    \toprule
    & \multicolumn{3}{c}{ABP-mean (mmHg) $\downarrow$} & \multicolumn{3}{c}{HR (bpm) $\downarrow$} \\
    \cmidrule(lr){2-4}\cmidrule(lr){5-7}
    Method & MCAR & MAR & NMAR & MCAR & MAR & NMAR \\
    \midrule
    \textbf{CAIR\ (Ours)} & \textbf{4.30} & \textbf{8.60} & \textbf{11.42} & \underline{2.26} & \textbf{5.28} & \underline{6.48} \\
    linear interp & \underline{4.40} & \underline{9.23} & \underline{11.85} & \textbf{2.22} & \underline{5.50} & \textbf{6.42} \\
    GP-VAE & 11.03 & 12.82 & 17.94 & 9.17 & 14.01 & 18.00 \\
    M-RNN & 13.45 & 13.19 & 18.35 & 15.28 & 14.53 & 18.60 \\
    \bottomrule
  \end{tabular}}
  \caption{\textbf{Cross-domain transfer to MIMIC-III ICU vitals}
  (reconstruction RMSE, mean over six missingness rates, lower is better).
  CAIR\ is the most accurate method on arterial pressure under every mechanism;
  heart rate is a boundary case. The CAIR\ row is the multivariate variant;
  both variants, the remaining baselines and per-mechanism Wilcoxon tests are
  in the supplement.
  \textbf{Bold} = best, \underline{underline} = second, per column.}
  \label{tab:mimic}
\end{table}

\begin{table}[htpb]
  \centering
  {\small
  \setlength{\tabcolsep}{3pt}
  \begin{tabular}{lccccr}
    \toprule
    Variant & MCAR & MAR & NMAR & \textbf{ALL} $\downarrow$ & $\Delta$ \\
    \midrule
    \textbf{CAIR\ (full)}    & \textbf{\aFullMcar} & \underline{\aFullMar} & \textbf{\aFullNmar} & \textbf{\aFullAll} & n/a \\
    \midrule
    \; w/o aux.\ loss         & \aAuxMcar & \aAuxMar & \aAuxNmar & \aAuxAll & \aDAux \\
    \; w/o residual           & \aResMcar & \aResMar & \aResNmar & \aResAll & \aDRes \\
    \; GRU $\to$ self-attn.   & \aAtnMcar & \textbf{\aAtnMar} & \underline{\aAtnNmar} & \underline{\aAtnAll} & \aDAtn \\
    \midrule
    linear interp             & \underline{\aLinMcar} & \aLinMar & \aLinNmar & \aLinAll & \aDLin \\
    \bottomrule
  \end{tabular}}
  \caption{\textbf{Component ablation} on AI-READI CGM, on the protocol of
  Table~\ref{tab:mech} (all $\aNpart$ test participants; RMSE mg/dL over six
  missingness rates). \textbf{ALL} is the mean of the three mechanisms,
  $\Delta$ the change against the full model. Rows are single training runs
  under the published recipe, differing only in the ablated component; linear
  is repeated from Table~\ref{tab:mech} for scale.
  \textbf{Bold} = best, \underline{underline} = second best, per metric column.}
  \label{tab:ablation}
\end{table}

\begin{table}[t]
  \centering
  {\small
  \setlength{\tabcolsep}{2.4pt}
  \begin{tabular}{lcccccc}
    \toprule
    Pretrain, then fine-tune & meal & sleep & asc & dip & comb & AVG \\
    \midrule
    control (AI-READI only)  & \textbf{15.60} & 27.52 & \textbf{5.96} & \textbf{5.14} & \textbf{14.74} & \textbf{13.79} \\
    $+$ HUPA-UCM (T1D)       & 15.99 & \textbf{26.32} & 6.21 & 5.36 & 15.66 & 13.91 \\
    $+$ OhioT1DM (T1D)       & 15.77 & \underline{27.14} & \underline{6.08} & 5.21 & 15.25 & \underline{13.89} \\
    $+$ Shanghai (T2D)       & \underline{15.71} & 27.58 & 6.13 & \underline{5.20} & \underline{15.18} & 13.96 \\
    \bottomrule
  \end{tabular}}
  \caption{\textbf{External-cohort pretraining does not improve accuracy}
  (physiological protocol, $352$ participants, single seed; RMSE mg/dL). Every
  external cell is within $0.2$\,mg/dL of the control and none improves on it.
  \textbf{Bold} = best, \underline{underline} = second best, per column.}
  \label{tab:pretrain}
\end{table}

\begin{table}[t]
  \centering
  {\small
  \setlength{\tabcolsep}{3.2pt}
  \begin{tabular}{lcccccc}
    \toprule
    Method & meal & sleep & asc & dip & comb & \textbf{AVG} $\downarrow$ \\
    \midrule
    \textbf{CAIR\ (Ours)} & \textbf{15.25} & \textbf{24.94} & \textbf{9.01} & \textbf{7.49} & \textbf{16.37} & \textbf{14.61} \\
    Savitzky-Golay & \underline{17.27} & \underline{29.45} & \underline{13.41} & \underline{11.22} & \underline{19.43} & \underline{18.15} \\
    \midrule
    EWMA          & 20.59 & 31.69 & 21.45 & 18.32 & 23.11 & 23.03 \\
    local mean    & 22.66 & 32.23 & 24.00 & 20.65 & 24.40 & 24.79 \\
    cubic spline  & 22.64 & 45.22 & 21.90 & 18.62 & 23.85 & 26.44 \\
    LOCF          & 25.91 & 37.33 & 34.19 & 29.36 & 31.81 & 31.72 \\
    \bottomrule
  \end{tabular}}
  \caption{\textbf{The four baselines omitted from Table~3 of the main
  paper} (AI-READI, physiological five-strategy protocol, all $352$ test
  participants; RMSE mg/dL, lower is better). The lower block is the omitted
  group; CAIR\ and Savitzky-Golay are repeated from Table~3 as the best method
  and the weakest one carried there. Every omitted method is at least
  $4.9$\,mg/dL behind that floor on the average. \textbf{Bold} = best,
  \underline{underline} = second best, per column.}
  \label{tab:phys-rest}
\end{table}

\begin{table}[t]
  \centering
  {\small
  \setlength{\tabcolsep}{4pt}
  \begin{tabular}{lccccc}
    \toprule
    Conditioning set & meal & sleep & asc & dip & \textbf{AVG} \\
    \midrule
    target only (control) & 14.94 & 25.17 & 6.07 & 4.93 & 13.11 \\
    $+$ heart rate        & 14.93 & 25.16 & 6.03 & 4.92 & 13.10 \\
    $+$ steps, calories   & 14.89 & 25.25 & 6.10 & 4.97 & 13.14 \\
    $+$ sleep state       & 14.89 & 25.27 & 5.98 & 4.92 & 13.09 \\
    $+$ respiration, stress & 14.79 & 24.96 & 6.00 & 4.89 & \textbf{12.98} \\
    $+$ environment       & 14.82 & 24.83 & 6.24 & 5.06 & 13.10 \\
    $+$ clinical          & 14.79 & 24.85 & 6.18 & 5.01 & \underline{13.06} \\
    $+$ ECG               & 14.86 & 24.87 & 6.26 & 5.09 & 13.14 \\
    $+$ retinal           & 14.86 & 24.90 & 6.12 & 4.98 & 13.07 \\
    \bottomrule
  \end{tabular}}
  \caption{Conditioning-modality ladder on AI-READI CGM (five-seed ensemble per
  rung; RMSE mg/dL). Adding modalities does not improve glucose imputation; the AVG spread ($12.98$-$13.14$) is within seed noise. AVG is over all five physiological strategies; the combined column is omitted for space. Absolute values are not comparable to Table~3 of the main paper. \textbf{Bold}/\underline{underline} mark the best/second AVG.}
  \label{tab:ladder}
\end{table}

\begin{table}[t]
  \centering
  {\small
  \setlength{\tabcolsep}{6pt}
  \begin{tabular}{lccc}
    \toprule
    Method & MCAR & MAR & NMAR \\
    \midrule
    CAIR\ (univariate)   & \textbf{4.26} & \textbf{8.58} & \underline{11.44} \\
    CAIR\ (multivariate) & \underline{4.30} & \underline{8.60} & \textbf{11.42} \\
    linear interp   & 4.40 & 9.23 & 11.85 \\
    missForest      & 4.62 & 9.60 & 12.91 \\
    MICE            & 4.88 & 9.02 & 12.65 \\
    $k$NN           & 5.27 & 8.91 & 12.69 \\
    hot-deck        & 5.82 & 9.11 & 12.81 \\
    LOCF            & 5.89 & 10.43 & 12.79 \\
    Fourier         & 5.94 & 11.28 & 13.56 \\
    mean            & 9.33 & 9.93 & 14.01 \\
    mode            & 11.42 & 11.82 & 15.11 \\
    GP-VAE          & 11.03 & 12.82 & 17.94 \\
    M-RNN           & 13.45 & 13.19 & 18.35 \\
    \bottomrule
  \end{tabular}}
  \caption{Complete MIMIC-III arterial-pressure reconstruction RMSE (mmHg, mean
  over six missingness rates). \textbf{Bold} = best, \underline{underline} = second best, per column.}
  \label{tab:mimic-full}
\end{table}

\begin{table}[t]
  \centering
  {\small
  \setlength{\tabcolsep}{6pt}
  \begin{tabular}{lccc}
    \toprule
    Method & MCAR & MAR & NMAR \\
    \midrule
    CAIR\ (univariate)   & \underline{2.26} & \underline{5.36} & 6.72 \\
    CAIR\ (multivariate) & \underline{2.26} & \textbf{5.28} & \underline{6.48} \\
    linear interp   & \textbf{2.22} & 5.50 & \textbf{6.42} \\
    missForest      & 2.37 & 7.40 & 9.25 \\
    MICE            & 2.49 & 7.55 & 9.63 \\
    $k$NN           & 2.69 & 7.09 & 9.21 \\
    hot-deck        & 2.90 & 7.19 & 9.28 \\
    LOCF            & 3.07 & 7.06 & 7.79 \\
    Fourier         & 3.31 & 7.17 & 7.86 \\
    mean            & 8.41 & 9.22 & 11.69 \\
    mode            & 10.88 & 10.39 & 14.10 \\
    GP-VAE          & 9.17 & 14.01 & 18.00 \\
    M-RNN           & 15.28 & 14.53 & 18.60 \\
    \bottomrule
  \end{tabular}}
  \caption{Complete MIMIC-III heart-rate reconstruction RMSE (bpm, mean over six
  missingness rates). The CAIR\ (multivariate) and linear rows are those of
  Table~6 of the main paper. CAIR\ is the most accurate method under MAR and
  linear interpolation under MCAR and NMAR, the two separated by at most
  $0.22$\,bpm; the per-mechanism significance tests for those differences are in
  Table~\ref{tab:mimic-wilcoxon}.
  \textbf{Bold} = best, \underline{underline} = second best, per column.}
  \label{tab:mimic-hr-full}
\end{table}

\begin{table}[t]
  \centering
  {\small
  \setlength{\tabcolsep}{6pt}
  \begin{tabular}{llcc}
    \toprule
    Signal & Mechanism & $\Delta$ & $p$ \\
    \midrule
    \multirow{3}{*}{ABP-mean (mmHg)}
      & MCAR & $-0.105$ & $7.9\times10^{-9}$ \\
      & MAR  & $-0.625$ & $4.1\times10^{-21}$ \\
      & NMAR & $-0.375$ & $2.9\times10^{-3}$ \\
    \midrule
    \multirow{3}{*}{Heart rate (bpm)}
      & MCAR & $+0.040$ & $7.0\times10^{-40}$ \\
      & MAR  & $-0.220$ & $0.83$ \\
      & NMAR & $+0.060$ & $0.37$ \\
    \bottomrule
  \end{tabular}}
  \caption{\textbf{Paired Wilcoxon signed-rank tests on MIMIC-III}, CAIR\
  (multivariate) versus linear interpolation, per (window, rate) pair.
  $\Delta$ is the mean of the per-pair RMSE differences in native units, which
  need not equal the difference of the aggregate RMSEs in Table~6 of the main
  paper because RMSE does not aggregate linearly; negative favors CAIR. All three arterial-pressure tests favor CAIR; on heart rate,
  linear interpolation is significantly better under MCAR and the remaining two
  differences are not significant ($\alpha=0.05$).}
  \label{tab:mimic-wilcoxon}
\end{table}

\begin{table}[t]
  \centering
  {\small
  \setlength{\tabcolsep}{4pt}
  \begin{tabular}{lcc}
  \toprule
  Method & RMSE (mmHg) $\downarrow$ & Burden MRR $\uparrow$ \\
  \midrule
  \multicolumn{3}{l}{\emph{Low RMSE, burden not preserved}} \\
  \quad linear interp & 11.85 & $+0.14$ \\
  \quad LOCF & 12.79 & $+0.07$ \\
  \quad Fourier & 13.56 & $+0.05$ \\
  \midrule
  \multicolumn{3}{l}{\emph{Recovers burden, high RMSE}} \\
  \quad MICE & 12.65 & $+0.63$ \\
  \quad $k$NN & 12.69 & $\mathbf{+0.65}$ \\
  \quad hot-deck & 12.81 & $\mathbf{+0.65}$ \\
  \quad GP-VAE & 17.94 & $+0.64$ \\
  \quad M-RNN & 18.35 & $+0.64$ \\
  \midrule
  \multicolumn{3}{l}{\emph{Both}} \\
  \quad CAIR\ (univariate) & \underline{11.44} & $\mathbf{+0.65}$ \\
  \quad \textbf{CAIR\ (multivar.)} & \textbf{11.42} & $+0.64$ \\
  \bottomrule
  \end{tabular}}
  \caption{\textbf{Only CAIR\ ranks among the best on both axes} (MIMIC-III ABP,
  NMAR, mean over six rates). Reconstruction RMSE (mmHg, lower better) and
  clinical-burden recovery (MRR on time in/above/below range, higher better).
  Interpolants match the RMSE of CAIR\ but fail to preserve the burden; tabular and neural
  imputers recover the burden at much higher RMSE, with $k$NN and hot-deck
  matching its recovery at ${\sim}11\%$ higher RMSE. \textbf{Bold} = best, \underline{underline} = second best, per column.}
  \label{tab:mimic-down}
\end{table}

\begin{table}[t]
  \centering
  {\small
  \setlength{\tabcolsep}{6pt}
  \begin{tabular}{lcc}
    \toprule
    Fill & MCAR AUROC & NMAR AUROC \\
    \midrule
    oracle (recorded trace) & 0.679 & 0.678 \\
    \midrule
    mean-fill      & 0.677 & 0.683 \\
    $k$NN          & 0.680 & 0.685 \\
    CAIR\ (multivar.) & 0.688 & 0.688 \\
    linear interp  & 0.692 & 0.688 \\
    PCHIP          & 0.692 & 0.689 \\
    \bottomrule
  \end{tabular}}
  \caption{\textbf{Hard outcomes do not reward reconstruction fidelity}
  (MIMIC-III, in-hospital mortality predicted from arterial pressure alone,
  $30\%$ missingness, 1D-CNN readout; $n=1200$ under MCAR and $1181$ under
  NMAR). The oracle is the recorded trace and is the ceiling for this task, yet
  four fills exceed it under MCAR and all five do under NMAR, because
  interpolation denoises the vital. No column is ordered as reconstruction
  accuracy would predict, and no best value is marked.}
  \label{tab:singlevital}
\end{table}

\begin{table}[t]
  \centering
  {\small
  \setlength{\tabcolsep}{4pt}
  \begin{tabular}{lcccc}
    \toprule
    & \multicolumn{2}{c}{RMSE (mg/dL) $\downarrow$} & \multicolumn{2}{c}{TIR recovery $\uparrow$} \\
    \cmidrule(lr){2-3}\cmidrule(lr){4-5}
    Study group & CAIR & linear & CAIR & linear \\
    \midrule
    healthy            & \textbf{18.29} & 21.49 & \textbf{0.165} & 0.097 \\
    pre-diabetes       & \textbf{15.90} & 21.15 & \textbf{0.467} & 0.382 \\
    oral medication    & \textbf{33.96} & 40.19 & \textbf{0.601} & 0.556 \\
    insulin-dependent  & \textbf{27.23} & 35.17 & \textbf{0.608} & 0.489 \\
    \bottomrule
  \end{tabular}}
  \caption{\textbf{The CGM result holds in every diabetes study group}
  (AI-READI, NMAR, all $352$ test participants). CAIR\ attains lower
  reconstruction error and higher time-in-range recovery than linear
  interpolation in all four groups. \textbf{Bold} = better of the two, per group
  and metric.}
  \label{tab:cohort}
\end{table}